% ============================================================================
% KAPITEL 2: Axiome der Informations-Weber-Theorie (IWT)
% ============================================================================

\chapter{Axiome der Informations-Weber-Theorie (IWT)}
\label{ch:axiome_iwt}

\section{Einleitung}
\label{sec:axiome_einleitung}

Die IWT ist eine fundamentale, diskrete Ur-Theorie. Sie postuliert keine Raumzeit, keine Materie, keine Felder – nur Information. Raum, Zeit, Materie und Energie emergieren aus der Struktur und Dynamik eines diskreten Informationsnetzwerks.

Die folgenden Axiome definieren die IWT vollständig. Sie sind logisch unabhängig und bilden die Grundlage für alle weiteren Ableitungen.

\section{Axiom 1: Diskretes universelles Informationsfeld}
\label{sec:axiom1}

Es existiert ein skalares, diskretes Informationsfeld

\[
I_k^{(n)},
\]

das die vollständige physikalische Realität beschreibt. Hierbei ist:

\begin{itemize}
    \item \( k = 1, \dots, N \) – Index der diskreten Informationsknoten,
    \item \( n \in \mathbb{N}_0 \) – diskreter Zeitindex (Update-Schritt),
    \item \( I_k^{(n)} \in \mathbb{R}^+ \) – dimensionsloser Informationswert.
\end{itemize}

Materie, Energie, Wellen, Geometrie und Dynamik sind Manifestationen der Struktur und Veränderung dieses Feldes.

\section{Axiom 2: Informationserhaltung}
\label{sec:axiom2}

Die Gesamtinformation ist erhalten:

\[
\sum_{k=1}^{N} I_k^{(n)} = \text{const} \quad \text{für alle } n.
\]

Information wird nicht erzeugt oder vernichtet – sie wird nur zwischen den Knoten umverteilt.

\section{Axiom 3: Informationsmetrik als emergente Geometrie}
\label{sec:axiom3}

Die physikalische Raumzeit ist nicht fundamental. Stattdessen entsteht eine effektive diskrete Metrik

\[
g_{kl}^{(n)}
\]

aus der lokalen und globalen Struktur des Informationsfeldes und seiner Kopplungsmatrix \( K_{kl}^{(n)} \). Die Metrik ist dynamisch und wird nicht vorgegeben.

\section{Axiom 4: Variationsprinzip der diskreten Informationsdynamik}
\label{sec:axiom4}

Die Dynamik des Informationsfeldes folgt aus einem universellen diskreten Informations-Lagrange-Funktional

\[
\mathcal{L}_d[I_k^{(n)}] = \mathcal{L}_{\text{local}} + \mathcal{L}_{\text{global}} + \mathcal{L}_{\text{fraktal}},
\]

mit drei fundamentalen Beiträgen:

\begin{itemize}
    \item \textbf{Lokaler Anteil (Weber-Struktur):}  
    \[
    \mathcal{L}_{\text{local}} = \frac{1}{2} \sum_{k,l} K_{kl}^{(n)} \Delta_{kl} I^{(n)} \Delta_{kl} I^{(n)},
    \]
    beschreibt direkte Informationsflüsse zwischen benachbarten Knoten.

    \item \textbf{Globaler Anteil (Bohm-Struktur):}  
    \[
    \mathcal{L}_{\text{global}} = -\frac{\lambda}{2} \sum_k \frac{\Delta^2 I_k^{(n)}}{I_k^{(n)}},
    \]
    beschreibt nicht-lokale Organisationsstrukturen über das gesamte Netzwerk.

    \item \textbf{Fraktaler Anteil (kosmische Skalierung):}  
    \[
    \mathcal{L}_{\text{fraktal}} = \mu \ln\!\bigl(1 + \gamma_{\text{eff}} G_{\text{eff}} L^2\bigr) \sum_k I_k^{(n)},
    \]
    beschreibt die skaleninvariante Struktur des Universums.
\end{itemize}

\section{Axiom 5: Dynamische Gleichung der Informationsmetrik}
\label{sec:axiom5}

Die Metrik entsteht aus der Variation des diskreten Funktionals. Die fundamentale Gleichung der Informationsmetrik in diskreter Form lautet:

\[
g_{kl}^{(n+1)} = g_{kl}^{(n)} + T \left[ \Delta_k I^{(n)} \Delta_l I^{(n)} - \lambda \frac{\Delta_{kl}^2 I^{(n)}}{I_{kl}^{(n)}} + \mu g_{kl}^{(n)} \ln\!\bigl(1 + \gamma_{\text{eff}} G_{\text{eff}} L^2\bigr) \right].
\]

Sie vereinigt lokale Weber-Dynamik, globale Bohm-Struktur und fraktale kosmische Skalierung.

\section{Axiom 6: Energie als emergenter Informationsfluss}
\label{sec:axiom6}

Energie ist keine fundamentale Größe. Sie entsteht als Erhaltungsgröße des diskreten Informationsflusses. Die diskrete Energie folgt aus der Zeitinvarianz des Informations-Lagrange-Funktionals:

\[
E^{(n)} = \sum_k \left[ \frac{\partial \mathcal{L}_d}{\partial (\Delta_t I_k^{(n)})} \Delta_t I_k^{(n)} - \mathcal{L}_d \right],
\]
und es gilt \( E^{(n+1)} = E^{(n)} \).

\section{Axiom 7: Naturkonstanten als emergente Skalierungsparameter}
\label{sec:axiom7}

Die fundamentalen Naturkonstanten sind keine Eingaben der Theorie, sondern feste Punkte der diskreten Informationsdynamik. Sie emergieren aus den Netzwerkparametern:

\begin{itemize}
    \item \( c \) – maximale Informationsflussrate,
    \item \( \hbar \) – globale Informationsgranularität,
    \item \( G \) – fraktale Kopplungsstärke,
    \item \( \alpha \) – Verhältnis lokaler zu globaler Kopplung,
    \item \( k_B \) – Informations-Temperatur-Skala.
\end{itemize}

Die expliziten Zusammenhänge werden in Kapitel~\ref{ch:naturkonstanten} angegeben.

\section{Axiom 8: Emergenz von Raum, Zeit und Dynamik}
\label{sec:axiom8}

Raum, Zeit und Dynamik sind emergente Eigenschaften der Informationsmetrik:

\begin{itemize}
    \item \textbf{Raum:} Der physikalische Raum entsteht aus der diskreten Metrik  
    \[
    d_{kl}^{(n)} = \sqrt{g_{kl}^{(n)}} \cdot \lambda_0,
    \]
    mit fundamentaler Länge \( \lambda_0 \).

    \item \textbf{Zeit:} Die physikalische Zeit entsteht als Ordnungsstruktur der Informationsänderung  
    \[
    t \approx n \cdot T,
    \]
    wobei \( n \) der Update-Index und \( T \) der fundamentale Zeitschrift ist.

    \item \textbf{Dynamik:} Klassische Mechanik, Quantenmechanik und Gravitation sind Grenzfälle der diskreten Informationsdynamik (siehe Anhang~\ref{app:vergleich}).
\end{itemize}

\section{Axiom 9: Universelle Gültigkeit}
\label{sec:axiom9}

Die IWT gilt auf allen Skalen:

\begin{itemize}
    \item \textbf{Mikroskopisch:} Quantenstruktur, Elementarteilchen, Wellenphänomene,
    \item \textbf{Mesoskopisch:} Klassische Mechanik, Elektrodynamik, Thermodynamik,
    \item \textbf{Makroskopisch:} Gravitation, Astrophysik, Planetenbewegung,
    \item \textbf{Kosmologisch:} Rotverschiebung, Hintergrundstrahlung, großskalige Struktur (siehe Kapitel~\ref{ch:kosmologie}).
\end{itemize}

\subsection{Vergleich mit den Axiomen der Standardphysik}
\label{sec:axiom_vergleich}

Die folgende Tabelle stellt die Axiome der Allgemeinen Relativitätstheorie (ART) 
denen der IWT/WDBT+ gegenüber. Wer die WDBT+ mit den Axiomen der ART misst, 
begeht einen Kategorienfehler.

\begin{tabular}{|l|l|}
\hline
\textbf{Standardaxiom (ART)} & \textbf{Axiom der IWT/WDBT+} \\
\hline
Raumzeit ist fundamental & Raumzeit emergiert aus Information \\
\hline
Gravitation = Krümmung & Gravitation = Weber-Kraft + \(\Omega\)-Feld (siehe Kapitel~\ref{ch:maxwell_gravitation}) \\
\hline
Lichtablenkung frequenzunabhängig & Lichtablenkung frequenzabhängig (\(\Delta\phi \propto \lambda^2\), siehe Kapitel~\ref{ch:testbare_vorhersagen}) \\
\hline
Expansion des Raumes & Stationäres Universum (keine Expansion, siehe Kapitel~\ref{ch:kosmologie}) \\
\hline
Dunkle Materie wird benötigt & Keine dunkle Materie (erklärt durch \(\Omega\)-Feld) \\
\hline
Dunkle Energie wird benötigt & Keine dunkle Energie (stationäre Kosmologie) \\
\hline
Zeitpfeil nicht erklärt (zeitumkehrinvariant) & Zeitpfeil durch \(\dot{\beta}>0\) (DSTT-Drift, siehe Kapitel~\ref{ch:dstt_weber_kraefte}) \\
\hline
Keine natürliche Grenze für kompakte Massen & Maximale Masse \(M_{\text{BEC}} \approx 3\times10^6 M_\odot\) (siehe Kapitel~\ref{ch:bec_objekte}) \\
\hline
Neutrinomasse nicht vorhergesagt & \(m_\nu \approx 0,1\,\text{eV}/c^2\) (siehe Kapitel~\ref{ch:neutrino}) \\
\hline
\end{tabular}

\vspace{0.3cm}
Die Frage ist nicht, ob die WDBT+ mit der ART übereinstimmt – das tut sie nicht, 
weil sie andere Axiome hat. Die Frage ist, ob sie funktioniert und ob ihre Vorhersagen 
bestätigt werden.

\section{Zusammenfassung}
\label{sec:axiome_zusammenfassung}

Die neun Axiome definieren die IWT vollständig in ihrer fundamentalen diskreten Formulierung. Alle physikalischen Phänomene – von der Quantenmechanik über die Gravitation bis zur Kosmologie – folgen aus der Struktur und Dynamik des diskreten Informationsfeldes und der daraus emergierenden Metrik.

Die IWT ist fundamental diskret, emergent kontinuierlich, vereinheitlicht alle Wechselwirkungen und macht testbare Vorhersagen (siehe Kapitel~\ref{ch:testbare_vorhersagen}).